\section{Limitations}

In this section, we summarize the limitations of the current implementation.
%
The limitations in this section have been detailed in other parts of this
document, but each was deemed significant enough to collect into this section
for additional visibility and quick reference.


\subsection{Geometry and Mesh}

\subsubsection{Collapsed Elements}

GFR does not support collapsed elements.
%
These are elements with one or more faces that have zero area.
%
A common situation where this occurs is when a quadrilateral planar
mesh is rotationally extruded about a singularity axis, often
referred to as a pole boundary.
% 
In many cases these collapsed elements can be replaced with other
supported polygon shapes with fewer faces.
%
However, this is best done within the mesh generation software
so that any curvature or high-order information is properly
included in the revised elements.


\subsection{Synthetic Eddy Method}

\subsubsection{Inflow Plane}

Any SEM inflow boundary must be on a constant $x$ plane.
%
Users should orient their grids to meet this restriction.


\subsubsection{Repeatability of Random Seed}

Because GFR uses the date\slash time to generate the initial seed,
it is not possible to reproduce the same sequence of random numbers
used by the synthetic eddy method.
%
This means that it is not possible to \emph{exactly} reproduce the same
results by rerunning a simulation from scratch.
%
However, over a long period, the time-averaged results should be in
very close agreement.
%
This would be similar to any differences one might see comparing time-averaged
results from one solution using two time-averaging windows of the same
duration but shifted in the averaging start time.


\devNotes{
For regression testing, we may want the ability to specify the initial seed.
Perhaps add an integer input:
\\
\begin{tabular}{ @{\qquad} lll}
    \inptext{Random\_Seed}
 &  $=0$ & Use date/time. \\
 &  $>0$ & Use specified value.
\end{tabular}
\\
and provide some guideline on how big a seed to specify
(if that affects the random entropy).
}


\subsection{Boundary Conditions}

\subsubsection{Periodic Boundaries}

Opposing periodic boundaries must have identical face elements in order to
have one-to-one connectivity between cell faces.



